\chapter{Hiện thực}
\label{chap:implementation}

\section{Môi trường và Công cụ Phát triển}
\label{sec:impl-env}

\begin{table}[H]
  \centering
  \caption{Công cụ và thư viện sử dụng trong dự án}
  \label{tab:impl-tools}
  \begin{tabular}{lll}
    \toprule
    \textbf{Công cụ / Thư viện} & \textbf{Phiên bản} & \textbf{Mục đích} \\
    \midrule
    React Native         & 0.81.5  & Framework ứng dụng di động \\
    Expo SDK             & 54.0.33 & Runtime, toolchain, build \\
    TypeScript           & 5.9.2   & Ngôn ngữ lập trình \\
    Expo Router          & 6.0.23  & File-based navigation \\
    NativeWind           & 4.2.3   & Styling (Tailwind CSS cho RN) \\
    Drizzle ORM          & 0.44.0  & ORM cho SQLite \\
    expo-sqlite          & --      & SQLite driver trên thiết bị \\
    Axios                & 1.17.0  & HTTP client giao tiếp backend \\
    AsyncStorage         & 2.2.0   & Lưu JWT token cục bộ \\
    Node.js Test Runner  & built-in & Unit test framework \\
    VS Code              & --      & IDE phát triển \\
    Git / GitHub         & --      & Quản lý mã nguồn \\
    \bottomrule
  \end{tabular}
\end{table}

\section{Cấu trúc Mã nguồn}
\label{sec:impl-structure}

Dự án được tổ chức theo cấu trúc \textit{feature-based},
tách biệt rõ ràng giữa UI, nghiệp vụ và dữ liệu:

\begin{verbatim}
src/
  app/                    # Màn hình (Expo Router file-based)
    _layout.tsx           # Root layout: AuthGuard, migration, seed
    (auth)/               # Luồng xác thực
      login.tsx, register.tsx
    (tabs)/               # Tab navigation chính
      index.tsx           # Trang chủ
      calc.tsx            # Module 1 - Tính chọn động cơ
      project.tsx         # Danh sách dự án
      profile.tsx         # Trang cá nhân
    chain.tsx             # Module 2 - Tính xích
    gear.tsx              # Module 3 - Tính bánh răng
    project/[id].tsx      # Chi tiết dự án
  lib/
    chainDriveCalculation.ts   # Nghiệp vụ tính xích (Facade + Strategy)
    gearDriveCalculation.ts    # Nghiệp vụ tính bánh răng
    globalState.ts             # Trạng thái liên module (Observer)
    token-storage.ts           # JWT persistence
  services/
    motor-dk-selection.ts      # Nghiệp vụ chọn động cơ (Module 1)
    api/                       # Axios services cho backend
  db/
    client.ts                  # Singleton Drizzle instance
    schema/motor_dk.ts         # Schema bảng motors_dk
    seed/                      # Dữ liệu động cơ ĐK và 4A
  context/
    auth-context.tsx           # Authentication provider
\end{verbatim}

\section{Hiện thực Module 1 -- Tính Chọn Động cơ}
\label{sec:impl-m1}

\subsection{Luồng Xử lý và Công thức Áp dụng}

Module~1 được hiện thực trong \texttt{src/services/motor-dk-selection.ts}
theo 6 bước của giáo trình~\cite{trinh2006tinh}:

\textbf{Bước 1 -- Tính hiệu suất truyền động tổng}:

Hệ dẫn động gồm: nối trục đàn hồi ($\eta_{nt} = 0{,}98$),
bánh răng côn ($\eta_{brc} = 0{,}96$), bánh răng trụ ($\eta_{brt} = 0{,}98$),
bộ truyền xích ($\eta_x = 0{,}92$), 3 cặp ổ lăn ($\eta_{ol} = 0{,}992$ mỗi cặp):
\[
  \eta = \eta_{nt} \cdot \eta_{brc} \cdot \eta_{brt} \cdot \eta_x \cdot \eta_{ol}^3
       = 0{,}98 \times 0{,}96 \times 0{,}98 \times 0{,}92 \times 0{,}992^3
       \approx 0{,}828
\]

\textbf{Bước 2 -- Công suất cần thiết}:
\[
  P_{ct} = \frac{P_{lv}}{\eta}
\]

\textbf{Bước 3 -- Số vòng quay sơ bộ}:
\[
  n_{sb} = u_{sb} \cdot n_{lv} = 28 \cdot n_{lv}
\]
với $u_{sb} = u_{hgt} \cdot u_x = 14 \times 2 = 28$.

\textbf{Bước 4 -- Chọn động cơ}: Hàm \texttt{selectMotorDkStep4Top()} lọc
$P_{đc} \geq P_{ct}$, nhóm theo (tốc độ đồng bộ, dòng sản phẩm),
lấy động cơ nhỏ công suất nhất mỗi nhóm, sắp xếp theo $|n_{đb} - n_{sb}|$ tăng dần.
Trả về tối đa 5 ứng viên để người dùng chọn.

\textbf{Bước 5 -- Tỉ số truyền thực tế}:
\[
  u_t = \frac{n_{dc}}{n_{lv}}, \quad
  u_h = \frac{u_t}{u_x}, \quad
  u_1 = \text{round}_{1\text{dp}}(0{,}24 \cdot u_h), \quad
  u_2 = \frac{u_h}{u_1}
\]

\textbf{Bước 6 -- Bảng $P$--$n$--$T$ trên 4 trục}:
Công suất được tính ngược từ trục công tác lên trục động cơ qua từng cặp hiệu suất;
số vòng quay theo chuỗi tỉ số truyền;
mô-men xoắn $T = 9{,}55 \times 10^6 \cdot P/n$ (N$\cdot$mm).

\subsection{Truy vấn CSDL Động cơ}

Hàm \texttt{selectMotorDkStep4Top} sử dụng Drizzle ORM để truy vấn SQLite:

\begin{verbatim}
const motors = await db.select().from(Motors_Dk)
  .where(gte(Motors_Dk.powerKw, P_ct));
\end{verbatim}

Sau đó áp dụng logic xếp hạng thuần TypeScript (không SQL) để nhóm theo
(syncSpeedRpm, series) và chọn ứng viên tốt nhất mỗi nhóm.

\section{Hiện thực Module 2 -- Tính Bộ truyền Xích}
\label{sec:impl-m2}

\subsection{Luồng Xử lý và Công thức Áp dụng}

Module~2 được hiện thực trong \texttt{src/lib/chainDriveCalculation.ts}
với hàm Facade \texttt{calculateChainDrive(input, constants)}.

Đầu vào: $P$ (kW), $n$ (v/p), $u$ từ kết quả Module~1 ($P_{III}$, $n_{III}$, $u_x$).

\textbf{Số răng đĩa xích}:
\[
  z_1 = \max(19,\; 29 - 2u), \quad z_2 = \min(120,\; \text{round}(u \cdot z_1))
\]

\textbf{Hệ số điều chỉnh công suất và bước xích tính toán}:
\[
  K_z = \frac{25}{z_1}, \quad
  K_n = \frac{n_{01}}{n}, \quad
  P_t = \frac{K \cdot K_z \cdot K_n \cdot P}{K_x}
\]
Tra bảng~5.4 (Bảng \texttt{TABLE\_5\_4} trong code)
tìm bước xích $p_c$ nhỏ nhất sao cho $P_{max}(n_{01}) \geq P_t$.

\textbf{Hình học bộ truyền}:
\[
  d_1 = \frac{p_c z_1}{\pi}, \quad
  d_2 = \frac{p_c z_2}{\pi}, \quad
  a_{sb} = 44 p_c
\]
\[
  X_{sb} = \frac{2a_{sb}}{p_c} + \frac{z_1+z_2}{2}
           + \left(\frac{z_2-z_1}{2\pi}\right)^2 \frac{p_c}{a_{sb}}
\]
Số mắt xích $X$ được làm tròn đến số chẵn.
Khoảng cách trục chính xác:
\[
  a = 0{,}25 p_c \left(X - M + \sqrt{(X-M)^2 - 8N}\right) - 0{,}002 \cdot a_{temp}
\]

\subsection{Hiện thực Kiểm nghiệm Độ bền}

Bốn điều kiện kiểm nghiệm, kết quả lưu trong \texttt{ChainDriveValidations}:

\begin{enumerate}
  \item \textbf{Va đập}: $i = \frac{z_1 n}{15X} \leq [i] = 20$ lần/s.
  \item \textbf{An toàn}: $s = \frac{Q_{sb}}{k_d F_t + F_v + F_o} \geq [s] = 8{,}2$.
  \item \textbf{Tiếp xúc}:
        $\sigma_H = 0{,}47\sqrt{\dfrac{k_r(k_d F_t + F_{vd}) \cdot E}{A \cdot K_d}}
         \leq [\sigma_H] = 600$ MPa.
  \item \textbf{Áp suất}: $p_{calc} \leq [p_0]$.
\end{enumerate}

\section{Hiện thực Module 3 -- Tính Bộ truyền Bánh răng}
\label{sec:impl-m3}

\subsection{Bánh răng Côn -- Cấp Nhanh}

Hiện thực trong hàm \texttt{calculateBevelGearStage()} với 7 bước:

\textbf{Bước 1 -- Vật liệu và ứng suất cho phép}:
Cả hai bánh dùng thép 40XH tôi cải thiện ($HB_1 = 300$, $HB_2 = 290$).
Số chu kỳ làm việc cơ sở:
$N_{HO} = 30 \cdot HB_2^{2{,}4}$; $N_{FO} = 4 \times 10^6$.
Ứng suất tiếp xúc cho phép:
\[
  [\sigma_H] = \frac{\sigma_{Hlim0}}{s_H} \cdot K_{HL}
             = \frac{2 \cdot 290 + 70}{1{,}1} \approx 590{,}9 \text{ MPa}
\]

\textbf{Bước 2 -- Kích thước sơ bộ}:
\[
  R_e = K_r \sqrt{u_1^2+1} \cdot
        \sqrt[3]{\frac{T_I K_{H\beta}}{(1-K_{be}) K_{be} u_1 [\sigma_H]^2}}
\]
với $K_{be} = 0{,}29$, $K_{H\beta} = 1{,}18$, $K_r = 50$, $K_d = 100$.

\textbf{Bước 3 -- Số răng và mô-đun}:
$z_{1p} = 16$; $z_1 = \text{round}_{even}(1{,}6 \cdot z_{1p}) = 26$;
$z_2 = \min(120, u_1 \cdot z_1)$.
Mô-đun vòng ngoài $m_{te}$ được chọn từ dãy tiêu chuẩn.

\textbf{Bước 4 -- Góc côn}: $\delta_1 = \arctan(z_1/z_2)$; $\delta_2 = 90° - \delta_1$.

\textbf{Bước 5--6 -- Kiểm nghiệm tiếp xúc và uốn}:
Tính $K_{Hv}$, $K_{Fv}$ (hệ số tải trọng động),
\[
  \sigma_H = Z_M Z_H Z_\varepsilon
             \sqrt{\frac{2 T_I K_H \sqrt{u_1^2+1}}{0{,}85 \cdot b \cdot d_{m1}^2 \cdot u_1}}
\]
\[
  \sigma_{F1} = \frac{2 T_I K_F Y_\varepsilon Y_\beta Y_{F1}}{0{,}85 \cdot b \cdot m_{nm} \cdot d_{m1}}
\]

\textbf{Bước 7 -- Lực tác dụng}:
$F_t = 2T_I/d_{m1}$; $F_n = F_t/\cos\alpha$;
$F_{r1} = F_t \tan\alpha \cos\delta_1$; $F_{a1} = F_t \tan\alpha \sin\delta_1$.

\subsection{Bánh răng Trụ -- Cấp Chậm}

Hiện thực trong hàm \texttt{calculateCylindricalGearStage()} với 5 bước chính.
Vật liệu: thép HB$_2 = 270$; $[\sigma_H] = (2\cdot270+70)/1{,}1 \cdot K_{HL}$.

\textbf{Khoảng cách trục sơ bộ}:
\[
  a_w = K_a (u_2+1) \sqrt[3]{\frac{T_{II} K_{H\beta}}{[\sigma_H]^2 u_2 \psi_{ba}}}
\]
với $K_a = 49{,}5$, $\psi_{ba} = 0{,}35$, $K_{H\beta} = 1{,}07$.
Làm tròn $a_w$ lên bội số của 5.

\textbf{Mô-đun}: $m \geq 0{,}01 a_w$, chọn giá trị tiêu chuẩn nhỏ nhất thỏa điều kiện
từ dãy $\{1,\; 1{,}25,\; 1{,}5,\; 2,\; 2{,}5,\; 3,\; 4,\; 5,\; 6,\; 8,\; 10\}$.

\section{Quản lý Trạng thái Liên Module}
\label{sec:impl-state}

File \texttt{src/lib/globalState.ts} định nghĩa kiểu \texttt{globalNavigationState}:

\begin{verbatim}
export const globalNavigationState: {
  gearResult: { bevelResult, cylindricalResult } | null;
  chainResult: ChainDriveResult | null;
  motorScreenState: { selectedMotor, u1, u2, ... } | null;
} = { gearResult: null, chainResult: null, motorScreenState: null };
\end{verbatim}

Đây là hiện thực của Observer Pattern:
khi Module~1 cập nhật state, Module~2 và 3 đọc và hiển thị kết quả
mà không cần truyền props qua stack navigation.

\section{Hiện thực Khởi động và Seed Dữ liệu}
\label{sec:impl-startup}

Root layout (\texttt{src/app/\_layout.tsx}) thực hiện chuỗi khởi tạo khi app mở:

\begin{enumerate}
  \item \texttt{migrate(db)} -- chạy SQL migration tạo bảng \texttt{motors\_dk} và \texttt{motors\_4a}.
  \item \texttt{seedMotorsDk(db)} -- kiểm tra nếu bảng trống thì nạp toàn bộ dữ liệu catalog.
  \item \texttt{seedMotors4a(db)} -- tương tự cho series 4A.
\end{enumerate}

Sau khi khởi tạo xong, splash screen được ẩn và \texttt{AuthGuard} kiểm tra token:
nếu chưa đăng nhập thì chuyển hướng đến \texttt{/(auth)/login}.
